\section{What the Identification Gap Costs}
\label{sec:theory}

This section states what we can prove about $\Delta^{\mathrm{id}}$ and, equally
important, where the proofs stop. Every claim carries its status: proved
unconditionally, proved under a stated and \emph{measured} constant, proved in a
special case that covers our benchmark, or open. We are explicit because two
earlier versions of this analysis overstated results of ours, and both corrections
are recorded below rather than removed.

\subsection{The geometric half, unconditional}

Proposition~\ref{prop:disjoint} of Section~\ref{sec:setup} gives the boundary:
with pairwise disjoint verbalizers, $\Delta^{\mathrm{id}}\equiv0$. Inside a shared
scope, Proposition~\ref{prop:cheb} of Section~\ref{sec:method} gives the exact
minimax quantity --- the Chebyshev radius $r_S$ of the member tasks' gauge-fixed
optima, attained uniquely at their Chebyshev centre, with
$r_S\ge\tfrac12\mathrm{diam}$ and $r_S=q_1(S)$. Both are unconditional, and
Proposition~\ref{prop:disjoint} is falsifiable: it predicts \emph{exactly} zero on
every singleton scope, and we check it in every run we report ($0$ violations,
$87$--$144$ observations per arm, worst absolute deviation $0.0$).

\subsection{The accuracy half needs a constant, and that constant is small}

The statement one wants next is that some task must \emph{lose accuracy} in
proportion to $r_S$. It is false without an assumption, and this is a correction to
an earlier form of our own claim, which bounded an accuracy difference below by a
diameter-type quantity. Accuracy differences are at most $1$; diameters are
unbounded. Concretely: two tasks with verbalizer margins $\pm10$ can be rescaled
to margins $\pm0.01$, shrinking every logit distance by $1000\times$ while the
accuracy loss stays exactly $1$. Distance and loss are different quantities, and no
assumption-free proportionality exists between them.

The repaired statement requires a margin density. Let the per-example verbalizer
margin have density at least $\kappa$ on $[0,\varepsilon]$.

\begin{proposition}[Accuracy-space minimax, conditional on $\kappa$]
\label{prop:cheb-acc}
Under the margin-density assumption, for any shared offset $\bb$ on scope $S$,
\begin{equation}
  \max_{g\in\mathrm{mem}(S)}\big[R_g^{\mathrm{orc}}-R_g(\bb)\big]
  \;\ge\; \kappa\cdot\min(r_S,\,t_0),
  \label{eq:cheb-acc}
\end{equation}
where $t_0$ is the largest shift for which the density bound holds.
\end{proposition}

We measured $\kappa$ rather than leaving it a symbol. On this stream the
worst-task value is $\kappa\approx\mathbf{0.0182}$ (WiC); the median over tasks is
$0.0625$, which \emph{overstates} the usable constant by $3.4\times$, because
$\kappa$ must hold simultaneously for every task sharing an offset and so the
smallest-$\kappa$ member sets the bound. Two further conservatisms: multi-class
flip fractions count only top-$2$ flips, and the estimator must take
$\min_i N(t_i)/t_{i+1}$ rather than $\min_i N(t_i)/t_i$, since
Proposition~\ref{prop:cheb-acc} needs the density at all $t$ and not only at grid
points --- the two differ by up to $2.5\times$ and only the former is a valid lower
bound.

\paragraph{The honest consequence.}
At $\kappa\approx0.018$, \eqref{eq:cheb-acc} is a weak bound on this benchmark. We
therefore do not lean on it: the geometric half carries the load, and every
accuracy-space number in this paper is \emph{measured} rather than derived from
\eqref{eq:cheb-acc}.

\subsection{Budgeted offsets: proved for binary shared verbalizers}

With $m$ centres per scope the relevant radius is the minimax $m$-centre covering
radius $q_m(S)$, and the ladder of Section~\ref{sec:ladder} is monotone in $m$ by
construction.

\begin{proposition}[Closed form and pigeonhole bound, $d=1$]
\label{prop:qm}
Let a scope have $K_S=2$ labels, so each gauge-fixed optimum is a scalar; sort them
$x_1\le\dots\le x_T$. Then $q_m$ equals the minimum over choices of $m-1$ cut
points of the largest resulting block half-range; for $m=2$,
$q_2=\min_j\max\!\big((x_j-x_1)/2,\ (x_T-x_{j+1})/2\big)$. A pigeonhole lower bound
on $q_m$ holds in every dimension and is \emph{tight} at $d=1$.
\end{proposition}

\noindent Status, stated precisely: proved for $d=1$ and any $m$; \textbf{open for
$m\ge2$ together with dimension $\ge2$}. We may write ``proved for binary shared
verbalizers'', never ``Proposition~\ref{prop:qm} is proved''. The pigeonhole bound
is strict rather than tight for $d\ge2$: $164$ of $400$ random instances disagree
with the enumerator, against $0$ of $600$ at $d=1$.

The special case suffices for every number we report, and for a reason that is a
property of the benchmark rather than a convenience: \emph{every scorable
multi-task scope on Order-4 has $K_S=2$}. The $K_S=3$ and $K_S=5$ scopes are
exactly the ones excluded from scoring --- CB's rarest class has $16$ examples, and
Yelp/Amazon's five-way verbalizers violate the distinct-first-piece precondition.
So every $q_m$ we report is backed by Proposition~\ref{prop:qm}, and the open case
is unmeasurable here rather than quietly assumed.

\paragraph{Two negative facts we pin with tests, because they narrow claims.}
First, the ``points $\cup$ pairwise midpoints'' candidate-centre enumerator is
exact only at $d=1$ or $n\le2$; at $d\ge2$ it \emph{overestimates} the radius (a
unit equilateral triangle gives $0.8660$ against the true $1/\sqrt3=0.5774$, a
factor of $1.5$). The direction favours us --- an overestimated radius inside a
lower bound means claiming less --- but it may not be called exact. Second, a
\emph{per-example} router can exceed the single-offset per-task oracle; we
construct a case going $0.500\to1.000$. So the accuracy-space corollary holds for
per-task routing only. That fact is also a live method question: whether a
task-free \emph{per-example} rule is achievable, we do not know.

\subsection{Why a rank budget cannot absorb this term}

Rank-based bounds on forgetting have the schematic form
\begin{equation}
  F(\theta)\;\le\;\|\nabla\ell\|\cdot\|\Delta\bW\|+\tfrac{L}{2}\|\Delta\bW\|^2,
  \qquad \|\Delta\bW\|\le\sqrt{\rho}\,\|\bB\|\|\bA\|,
  \label{eq:rankbound}
\end{equation}
and are the basis on which allocators such as E$^2$-LoRA, AdaLoRA and CoDyRA spend
rank. We previously asserted a strict impossibility here --- that
$\Delta^{\mathrm{id}}$ can be made large with $\|\Delta\bW\|$ arbitrarily small ---
and that is \textbf{not correct as stated}. A small $\|\Delta\bW\|$ does imply a
small logit perturbation, which does imply a small accuracy change \emph{unless}
many examples sit within that perturbation of the decision boundary. The correct
statement is conditional on the same $\kappa$:

\begin{proposition}[Margin-conditional amplification]
\label{prop:amp}
With margin density at least $\kappa$ on $[0,\varepsilon]$, a logit shift of size
$\delta\le\varepsilon$ along the verbalizer-difference direction flips at least
$\kappa\delta$ mass, and $\delta\le c\,\|\Delta\bW\|\sup_x\|\bh(x)\|$.
\end{proposition}

So \eqref{eq:rankbound} is loose by the factor $\kappa$: a rank-$1$ update aligned
with a shared verbalizer-difference direction can satisfy the rank bound and still
move accuracy by $\kappa\delta$. The claim we are entitled to is therefore weaker
than impossibility and sharper than vacuity: the rank bound is
\textbf{uninformative about which of the two terms the damage lands in}, and can be
satisfied while $\Delta^{\mathrm{id}}$ is large.

That is the operative point for allocators. Output-drift energy, importance scores
and rank minimisation are all averages over the whole drift; none separates the
component a zero-rank shared offset could absorb from the component that genuinely
needs subspace capacity. Under a fixed lifetime budget, rank spent on the former is
wasted --- and, by Propositions~\ref{prop:cheb} and~\ref{prop:qm}, offset capacity
spent on conflicting tasks cannot substitute for rank either. Both budgets bind,
for different reasons, which is the two-constraint picture of
Section~\ref{sec:setup}.